\chapter{Théorie cinétique et gaz parfait}\label{ch:b1:kinetic-theory}

L'air d'une salle de classe pèse autant qu'un homme adulte et contient
plus de molécules qu'il n'y a de grains de sable sur Terre, chacune
filant à la vitesse d'une balle de fusil et heurtant ses voisines dix
milliards de fois par seconde. Personne ne peut en suivre une~; personne
n'en a besoin. Une poignée de moyennes --- \omterm{def:b1:kinetic-theory:pressure}{pression}, \omterm{def:b1:kinetic-theory:temperature}{température},
volume, quantité de matière --- décrit l'air complètement pour tous les
usages de l'ingénieur, et une image simple de molécules qui rebondissent
sur les parois explique d'où viennent ces moyennes et comment elles sont
reliées. Ce chapitre ouvre la partie thermodynamique du volume en
regardant la matière par les deux bouts~: le chaos microscopique et le
calme macroscopique, et la théorie cinétique qui les relie dans le cas
du \omterm{def:b1:kinetic-theory:model}{gaz parfait}.

\section{Deux échelles de description}

\begin{definition}[Échelles microscopique et macroscopique]\label{def:b1:kinetic-theory:scales}
\index{échelle microscopique}\index{échelle macroscopique}\index{constante d'Avogadro}\index{mole}
À l'échelle \emph{microscopique}, la matière est faite de molécules~:
$\qty{1}{mol}$ en contient $N_A = \qty{6.02e23}{}$, chacune large
d'environ \qty{0.3}{nm}~; dans un gaz aux conditions ordinaires, elles
sont distantes de \qty{3}{nm}, se déplacent à des centaines de mètres
par seconde et se heurtent tous les \qty{70}{nm} environ. À l'échelle
\emph{macroscopique}, un système est décrit par quelques
\emph{\omterm{def:b1:kinetic-theory:equilibrium}{variables d'état}} --- \omterm{def:b1:kinetic-theory:pressure}{pression} $P$, volume $V$, \omterm{def:b1:kinetic-theory:temperature}{température} $T$,
quantité de matière $n$ --- qui sont des moyennes sur des nombres
énormes de molécules. Une variable est \emph{extensive} si elle double
quand le système double ($V$, $n$, masse, énergie), \emph{intensive}
sinon ($P$, $T$, masse volumique).
\end{definition}

\begin{definition}[Équilibre thermodynamique]\label{def:b1:kinetic-theory:equilibrium}
\index{équilibre thermodynamique}\index{variable d'état}
Un système est en \emph{équilibre thermodynamique} lorsque ses variables
d'état sont uniformes et constantes dans le temps et qu'aucun écoulement
macroscopique de matière ou d'énergie ne le traverse. C'est seulement
alors que $P$ et $T$ ont un sens pour le système entier~; l'\emph{équation
d'état} les relie à $V$ et à $n$.
\end{definition}

\begin{definition}[Pression]\label{def:b1:kinetic-theory:pressure}
\index{pression}\index{pascal}\index{bar}
Un fluide au repos pousse sur chaque élément de surface $\dd S$ d'une
paroi avec une force $\dd\vect F = P\,\dd S\,\vect n$ normale à celle-ci
et dirigée vers elle~; $P$ est la \emph{pression}, en pascals
($\qty{1}{Pa} = \qty{1}{N/m^2}$)~; $\qty{1}{bar}
= \qty{1e5}{Pa}$, $\qty{1}{atm} = \qty{1.013e5}{Pa}$. En un point donné,
elle ne dépend pas de l'orientation de la surface (\cref{ch:b1:fluid-statics}).
\end{definition}

\begin{definition}[Température]\label{def:b1:kinetic-theory:temperature}
\index{température}\index{kelvin}\index{principe zéro}
Deux systèmes en contact par une paroi qui laisse passer l'énergie
atteignent, à la longue, un état commun~: ils sont alors à la même
\emph{température}, et deux corps chacun en équilibre avec un troisième
sont en équilibre entre eux (le \emph{principe zéro}). L'échelle des
kelvins est fixée par $k_B = \qty{1.380649e-23}{J/K}$ exactement~; en
degrés Celsius, $\theta = T - 273.15$. La théorie cinétique qui suit
donne à $T$ son sens microscopique.
\end{definition}

\section{Le modèle cinétique du gaz parfait}

\begin{definition}[Gaz parfait, modèle microscopique]\label{def:b1:kinetic-theory:model}
\index{gaz parfait}\index{gaz idéal}
Un \emph{gaz parfait} est un ensemble de $N$ molécules, traitées comme
des points de masse $m$, se déplaçant librement et au hasard --- sans
interaction hormis de brefs chocs --- dans un récipient de volume $V$. À
l'équilibre, la distribution des vitesses est isotrope et stationnaire~;
$n^* = N/V$ est la densité particulaire et $\langle v^2\rangle$ la
\omterm{def:b1:kinetic-theory:kinetictemp}{vitesse quadratique moyenne}.
\end{definition}

\begin{theorem}[Pression cinétique]\label{thm:b1:kinetic-theory:pressure}
\index{pression cinétique}
La \omterm{def:b1:kinetic-theory:pressure}{pression} qu'exerce le gaz sur les parois vaut
\[
P = \tfrac13\,n^*m\langle v^2\rangle = \tfrac23\,n^*\langle\epsilon_k\rangle ,
\]
où $\langle\epsilon_k\rangle = \tfrac12 m\langle v^2\rangle$ est l'\omterm{thm:b1:work-and-energy:ke}{énergie
cinétique} moyenne d'une molécule.
\end{theorem}

\begin{proof}[Démonstration partielle]
Prenons le modèle simplifié où toutes les molécules ont la vitesse $u$
et se déplacent selon les trois axes, un sixième dans chaque sens. Une
molécule qui frappe de plein fouet la paroi $x = 0$ rebondit
élastiquement et lui cède la quantité de mouvement $2mu$. Pendant $\dd t$,
les molécules qui vont vers la paroi et atteignent une aire $S$ sont
celles de la tranche d'épaisseur $u\,\dd t$, soit un sixième des
$n^*Su\,\dd t$ qui s'y trouvent. Force $= $
quantité de mouvement par unité de temps $= (n^*Su\,\dd t/6)(2mu)/\dd t = \tfrac13 n^*mu^2S$, d'où
$P = \tfrac13 n^*mu^2$. Moyenner sur la vraie distribution des vitesses
et des directions remplace $u^2$ par $\langle v^2\rangle$ (l'isotropie
donne $\langle v_x^2\rangle = \tfrac13\langle v^2\rangle$)~; le volume de
deuxième année le fait complètement.
\end{proof}

\begin{omfigure}
\begin{tikzpicture}[font=\small, x=1cm, y=1cm]
  \draw[thick] (0,-1.6) -- (0,1.6); \fill[pattern=north east lines] (-0.2,-1.6) rectangle (0,1.6);
  \draw[dashed] (1.8,-1.6) -- (1.8,1.6);
  \draw[Stealth-Stealth] (0,-1.85) -- (1.8,-1.85) node[midway, below] {$u\,\dd t$};
  \foreach \p in {(0.6,1.1),(1.2,0.5),(0.9,-0.2),(1.5,-0.9),(0.4,-1.2),(1.3,1.3)} {\fill[omDef] \p circle (2pt);}
  \draw[-Stealth, omDef, thick] (1.2,0.5) -- (0.5,0.5);
  \draw[-Stealth, omThm, thick] (0.25,0.2) -- (0.95,0.2);
  \node[omDef, font=\footnotesize, above] at (0.85,0.55) {$m\vect u$};
  \node[omThm, font=\footnotesize, below] at (0.6,0.15) {$-m\vect u$};
  \node[right, font=\footnotesize] at (2.0,0.9) {quantité de mouvement $2mu$ cédée};
  \node[right, font=\footnotesize] at (2.0,0.55) {à la paroi par choc};
  \node[right, font=\footnotesize] at (2.0,-0.5) {$\tfrac16 n^*Su\,\dd t$ chocs pendant $\dd t$};
  \begin{scope}[xshift=7.2cm, yshift=-1.6cm]
  \begin{axis}[omaxis, axis x line=bottom, axis y line=left, width=6.4cm, height=4.2cm, xlabel={$v$ (\unit{m/s})}, ylabel={$f(v)$},
      xlabel style={at={(axis description cs:0.5,-0.1)}, anchor=north},
      xmin=0, xmax=1400, ymin=0, ymax=1.2, xtick={500,1000}, ytick=\empty, anchor=origin, at={(0,0)}]
    \addplot[omDef, thick, domain=0:1400, samples=120] {(x/422)^2*exp(1-(x/422)^2)};
    \addplot[omThm, thick, domain=0:1400, samples=120] {0.707*(x/597)^2*exp(1-(x/597)^2)};
    \node[omDef, font=\footnotesize] at (axis cs:430,1.12) {\qty{300}{K}};
    \node[omThm, font=\footnotesize] at (axis cs:900,0.72) {\qty{600}{K}};
  \end{axis}
  \end{scope}
\end{tikzpicture}

{\small À gauche~: l'origine cinétique de la \omterm{def:b1:kinetic-theory:pressure}{pression} --- les molécules
qui frappent une paroi et rebondissent lui cèdent de la quantité de
mouvement~; leur nombre par seconde multiplié par $2mu$ est la force. À
droite~: les vitesses réelles sont réparties (distribution de Maxwell,
ici pour l'azote à deux \omterm{def:b1:kinetic-theory:temperature}{températures})~; l'unique $u$ du modèle est
remplacé par la \omterm{def:b1:kinetic-theory:kinetictemp}{vitesse quadratique moyenne}.}
\end{omfigure}

\begin{definition}[Température cinétique~; équation d'état]\label{def:b1:kinetic-theory:kinetictemp}
\index{température cinétique}\index{équation d'état!du gaz parfait}\index{constante des gaz parfaits}\index{vitesse quadratique moyenne}
La \omterm{def:b1:kinetic-theory:temperature}{température} d'un \omterm{def:b1:kinetic-theory:model}{gaz parfait} est \emph{définie} par l'\omterm{thm:b1:work-and-energy:ke}{énergie
cinétique} moyenne de ses molécules,
\[
\langle\epsilon_k\rangle = \tfrac12 m\langle v^2\rangle = \tfrac32 k_BT ,
\]
de sorte que, en combinant avec le \cref{thm:b1:kinetic-theory:pressure},
\[
PV = Nk_BT = nRT, \qquad R = N_Ak_B = \qty{8.314}{J/(mol.K)} ,
\]
c'est l'\emph{équation d'état du \omterm{def:b1:kinetic-theory:model}{gaz parfait}}. La \emph{vitesse
quadratique moyenne} vaut $u = \sqrt{\langle v^2\rangle} = \sqrt{3k_BT/m} = \sqrt{3RT/M}$,
où $M$ est la masse molaire.
\end{definition}

\begin{example}[Quelques nombres pour l'air]\label{ex:b1:kinetic-theory:air}
Azote à \qty{300}{K}~: $u = \sqrt{3 \times 8.314 \times 300/0.028} =
\qty{517}{m/s}$~; hydrogène, $\qty{1930}{m/s}$~; plus c'est lourd, plus
c'est lent, en $1/\sqrt M$. À \qty{1}{\omterm{def:b1:kinetic-theory:pressure}{bar}} et \qty{300}{K}, $n^* = P/k_BT = \qty{2.4e25}{m^{-3}}$~:
\qty{24}{L} par \omterm{def:b1:kinetic-theory:scales}{mole}, \qty{2.7e22}{} molécules dans un litre. Une salle
de \qty{200}{m^3}~: \qty{8300}{mol}, soit \qty{240}{kg} d'air.
\end{example}

\begin{remark}[Pourquoi le modèle marche, et quand il échoue]\label{rem:b1:kinetic-theory:validity}
\index{libre parcours moyen}
Entre deux chocs, une molécule d'air vole sur un \emph{\omterm{rem:b1:kinetic-theory:validity}{libre parcours
moyen}} $\lambda \approx 1/(\sqrt2\pi d^2n^*) \approx \qty{70}{nm}$ ---
deux cents fois sa taille~: le gaz est surtout du vide, les interactions
sont rares, et $PV = nRT$ vaut à mieux qu'un pour cent aux \omterm{def:b1:kinetic-theory:pressure}{pressions}
ordinaires. À haute \omterm{def:b1:kinetic-theory:pressure}{pression} ou à basse \omterm{def:b1:kinetic-theory:temperature}{température}, le volume propre
des molécules et leur attraction comptent~: l'équation de van der Waals
$(P + an^2/V^2)(V - nb) = nRT$ corrige les deux, et sous une \omterm{def:b1:kinetic-theory:temperature}{température}
critique le gaz se liquéfie (\cref{ch:b1:phase-changes}).
\end{remark}

\begin{omfigure}
\begin{tikzpicture}
\begin{axis}[omaxis, axis x line=bottom, axis y line=left, width=10cm, height=6cm, xlabel={$V$}, ylabel={$P$},
    xlabel style={at={(axis description cs:0.5,-0.08)}, anchor=north},
    xmin=0, xmax=5.2, ymin=-0.4, ymax=4.2, xtick=\empty, ytick=\empty]
  \addplot[omDef, thick, domain=0.3:5.2, samples=100] {1.0/x};
  \addplot[omDef, thick, domain=0.45:5.2, samples=100] {1.5/x};
  \addplot[omDef, thick, domain=0.6:5.2, samples=100] {2.0/x};
  \addplot[omThm, thick, domain=0.75:5.2, samples=100] {2.5/x};
  \node[omDef, font=\footnotesize, below] at (axis cs:4.9,0.19) {$T_1$};
  \node[omThm, font=\footnotesize, above] at (axis cs:4.9,0.53) {$T_4$};
  \draw[-Stealth, omExo] (axis cs:1.3,1.0) -- (axis cs:2.3,2.0) node[above right, font=\footnotesize] {$T$ croissante};
  \node[font=\footnotesize] at (axis cs:3.4,3.2) {$PV = nRT$~: isothermes $T_1 < \dots < T_4$};
\end{axis}
\end{tikzpicture}

{\small Le diagramme de Clapeyron d'un \omterm{def:b1:kinetic-theory:model}{gaz parfait}~: à chaque
\omterm{def:b1:kinetic-theory:temperature}{température}, l'isotherme est une hyperbole $P = nRT/V$, d'autant plus
haute que $T$ est grande. Une transformation est un chemin dans ce plan~;
un cycle, une boucle fermée (\cref{ch:b1:heat-engines}).}
\end{omfigure}

\section{Énergie interne du gaz parfait}

\begin{definition}[Énergie interne]\label{def:b1:kinetic-theory:U}
\index{énergie interne}
L'\emph{énergie interne} $U$ d'un système est la somme des \omterm{thm:b1:work-and-energy:ke}{énergies
cinétiques} de ses molécules dans le \omterm{def:b1:point-kinematics:frame}{référentiel} où le système est au
repos et de leurs \omterm{def:b1:work-and-energy:conservative}{énergies potentielles} mutuelles~; c'est une fonction
d'état extensive.
\end{definition}

\begin{proposition}[Énergie interne des gaz parfaits]\label{prop:b1:kinetic-theory:Ugas}
\index{gaz monoatomique}\index{gaz diatomique}\index{capacité thermique à volume constant}\index{lois de Joule}
Pour un \omterm{def:b1:kinetic-theory:model}{gaz parfait}, $U$ ne dépend que de $T$ (première loi de Joule)~:
\begin{gather*}
U = \tfrac32 nRT \quad (\text{monoatomique~: He, Ne, Ar}), \\
U = \tfrac52 nRT \quad (\text{diatomique aux températures ordinaires~: N}_2, \text{O}_2, \text{H}_2),
\end{gather*}
si bien que la capacité thermique molaire à volume constant
$C_{V,m} = \dd U_m/\dd T$ vaut respectivement
$\tfrac32 R = \qty{12.5}{J/(mol.K)}$ et $\tfrac52 R = \qty{20.8}{J/(mol.K)}$.
\end{proposition}

\begin{proof}[Démonstration partielle]
Aucune \omterm{def:b1:work-and-energy:conservative}{énergie potentielle} mutuelle par hypothèse, donc $U = N\langle\epsilon_k\rangle
= \tfrac32 Nk_BT$ pour des molécules ponctuelles. Une molécule diatomique
tourne aussi~: chacun de ses deux degrés de liberté de rotation porte, à
l'équilibre, le même $\tfrac12 k_BT$ que chaque degré de translation
(théorème de l'\emph{équipartition}, admis~; le volume de troisième année
l'établit et explique pourquoi la vibration ne s'y joint qu'au-dessus du
millier de \omterm{def:b1:kinetic-theory:temperature}{kelvins}).
\end{proof}

\begin{example}[Ce que contient une salle]\label{ex:b1:kinetic-theory:room}
Les \qty{8300}{mol} d'air de la salle de classe à \qty{300}{K}~: $U = \tfrac52 \times
8300 \times 8.314 \times 300 = \qty{5.2e7}{J}$ --- l'énergie d'un litre et
demi d'essence, tenue dans l'agitation moléculaire. Réchauffer la salle
de $1$ K coûte $\tfrac52 nR = \qty{170}{kJ}$ (à volume constant).
\end{example}

\section{Les phases condensées}

\begin{proposition}[Modèle de la phase incompressible et indilatable]\label{prop:b1:kinetic-theory:condensed}
\index{phase condensée}\index{dilatation thermique}\index{compressibilité}
Les liquides et les solides sont modélisés, à ce niveau, par des phases
de volume fixe --- ni compressibles ni dilatables --- dont l'\omterm{def:b1:kinetic-theory:U}{énergie
interne} ne dépend que de $T$~: $\dd U = C\,\dd T$ avec $C = mc$, où $c$
est la capacité thermique massique (\qty{4180}{J/(kg.K)} pour l'eau,
\qty{450}{J/(kg.K)} pour l'acier). Les petits écarts réels se mesurent
par le coefficient de \omterm{prop:b1:kinetic-theory:condensed}{dilatation thermique}
$\alpha = (1/V)(\partial V/\partial T)_P$ ($\sim\qty{1e-5}{K^{-1}}$ pour
les métaux, \qty{2e-4}{K^{-1}} pour les liquides) et par la
compressibilité isotherme $\chi_T = -(1/V)(\partial V/\partial P)_T$
($\sim\qty{5e-10}{Pa^{-1}}$ pour l'eau, contre
$1/P \approx \qty{1e-5}{Pa^{-1}}$ pour un gaz).
\end{proposition}
\admitted

\begin{example}[À quel point incompressible]\label{ex:b1:kinetic-theory:water}
L'eau au fond de la fosse des Mariannes (\qty{1100}{\omterm{def:b1:kinetic-theory:pressure}{bar}}) est comprimée
de $\chi_T\Delta P \approx 5\times10^{-10} \times 1.1\times10^8 = 5\%$~; un
rail d'acier de \qty{30}{m} réchauffé de $50$ K s'allonge de $\alpha L\Delta T =
\qty{18}{mm}$ --- d'où les joints de dilatation. Face à un gaz, dont le
volume est divisé par deux quand la \omterm{def:b1:kinetic-theory:pressure}{pression} double, le modèle du volume
fixe est bon à quelques pour cent près dans presque toutes les
situations de ce volume.
\end{example}

\section{Exercices}

\begin{exercise}[$\star$]\label{exo:b1:kinetic-theory:1}
Nombre de molécules dans \qty{1.0}{L} de gaz à \qty{0}{\degreeCelsius} et
\qty{1}{atm}~; distance moyenne entre voisines~; comparer à la taille
moléculaire \qty{0.3}{nm}.
\end{exercise}

\begin{exercise}[$\star$]\label{exo:b1:kinetic-theory:2}
\omterm{def:b1:kinetic-theory:kinetictemp}{Vitesses quadratiques moyennes} de H$_2$, N$_2$, O$_2$ et CO$_2$ à
\qty{300}{K}, et de N$_2$ à \qty{77}{K} (azote bouillant).
\end{exercise}

\begin{exercise}[$\star$]\label{exo:b1:kinetic-theory:3}
Un pneu de voiture de \qty{30}{L} est gonflé à une \omterm{def:b1:kinetic-theory:pressure}{pression} relative de
\qty{2.5}{\omterm{def:b1:kinetic-theory:pressure}{bar}} à \qty{20}{\degreeCelsius}. Quantité de matière et masse
d'air à l'intérieur~; \omterm{def:b1:kinetic-theory:pressure}{pression} absolue après un trajet qui la porte à
\qty{50}{\degreeCelsius}.
\end{exercise}

\begin{exercise}[$\star$]\label{exo:b1:kinetic-theory:4}
Estimer le \omterm{rem:b1:kinetic-theory:validity}{libre parcours moyen} des molécules d'air à \qty{1}{\omterm{def:b1:kinetic-theory:pressure}{bar}} et
\qty{300}{K} ($d = \qty{0.37}{nm}$), ainsi que la fréquence des chocs
d'une molécule. À quelle \omterm{def:b1:kinetic-theory:pressure}{pression} le \omterm{rem:b1:kinetic-theory:validity}{libre parcours moyen} atteint-il
\qty{1}{m} (un «~vide~»)~?
\end{exercise}

\begin{exercise}[$\star\star$]\label{exo:b1:kinetic-theory:5}
Masse volumique de l'air à \qty{1.013}{\omterm{def:b1:kinetic-theory:pressure}{bar}} et \qty{20}{\degreeCelsius}
($M = \qty{29}{g/mol}$), puis à \qty{100}{\degreeCelsius}. Force
ascensionnelle d'une montgolfière de \qty{2800}{m^3} (différence des
deux poids d'air).
\end{exercise}

\begin{exercise}[$\star\star$]\label{exo:b1:kinetic-theory:6}
L'air contient $21\%$ de O$_2$ en molécules. \omterm{def:b1:kinetic-theory:pressure}{Pression} partielle de
l'oxygène au niveau de la mer~; à \qty{5000}{m} où $P = \qty{0.54}{atm}$~;
au sommet de l'Everest (\qty{0.33}{atm}). Pourquoi le corps y peine-t-il~?
\end{exercise}

\begin{exercise}[$\star\star$]\label{exo:b1:kinetic-theory:7}
Une \omterm{def:b1:kinetic-theory:scales}{mole} de CO$_2$ à \qty{300}{K} dans \qty{1.0}{L}~: \omterm{def:b1:kinetic-theory:pressure}{pression} par la loi
du \omterm{def:b1:kinetic-theory:model}{gaz parfait}, puis par van der Waals ($a = \qty{0.364}{Pa.m^6/mol^2}$,
$b = \qty{4.27e-5}{m^3/mol}$). Quelle correction domine ici~?
\end{exercise}

\begin{exercise}[$\star\star$]\label{exo:b1:kinetic-theory:8}
\omterm{def:b1:kinetic-theory:U}{Énergie interne} de l'air d'une pièce de \qty{50}{m^3} à \qty{1}{\omterm{def:b1:kinetic-theory:pressure}{bar}} et
\qty{300}{K}~; énergie pour la réchauffer de \qty{5}{K} à volume
constant~; durée avec un radiateur de \qty{2}{kW} (sans pertes).
\end{exercise}

\begin{exercise}[$\star\star$]\label{exo:b1:kinetic-theory:9}
Un rail d'acier de \qty{30}{m} ($\alpha = \qty{1.2e-5}{K^{-1}}$) entre
\qty{-10}{\degreeCelsius} et \qty{40}{\degreeCelsius}~: variation de
longueur. Une bague de laiton ($\alpha = \qty{1.9e-5}{K^{-1}}$) de
diamètre intérieur \qty{49.95}{mm} doit se glisser sur un arbre de
\qty{50.00}{mm}~: de combien faut-il la chauffer~?
\end{exercise}

\begin{exercise}[$\star\star\star$]\label{exo:b1:kinetic-theory:10}
Énergie par molécule à \qty{300}{K} ($\tfrac32 k_BT$) en joules et en eV~;
capacités thermiques molaires $C_{V,m}$ de l'argon et de l'azote~;
pourquoi celle de l'azote est-elle plus grande alors que ses molécules
sont plus légères~?
\end{exercise}

\begin{exercise}[$\star\star\star$]\label{exo:b1:kinetic-theory:11}
Refaire l'établissement de la \omterm{thm:b1:kinetic-theory:pressure}{pression cinétique} dans le modèle à six
directions, puis montrer que pour une distribution isotrope des vitesses $\langle v_x^2\rangle
= \tfrac13\langle v^2\rangle$ et que la formule de la \omterm{def:b1:kinetic-theory:pressure}{pression} est inchangée.
\end{exercise}

\begin{exercise}[$\star\star\star$]\label{exo:b1:kinetic-theory:12}
\omterm{def:b1:kinetic-theory:kinetictemp}{Vitesses quadratiques moyennes} de H$_2$ et de O$_2$ à \qty{300}{K},
comparées à la \omterm{prop:b1:central-forces:escape}{vitesse de libération} de la Terre puis à celle de la Lune
(\qty{2.4}{km/s}). Sachant que la fraction de molécules plus rapides que
$3u$ est petite mais non nulle, expliquer pourquoi la Terre a perdu son
hydrogène et la Lune toute son atmosphère.
\end{exercise}

\section{Problème~: l'air de la salle}

\begin{problem}[{Problème du week-end --- une salle de classe pleine
d'air~: compter ses molécules, la peser, chronométrer leurs vols, et
retrouver du même modèle la portance d'une montgolfière, la rareté de
l'air en altitude et la vitesse du son}]
\label{pb:b1:kinetic-theory:1}
Une salle de classe mesure $10 \times 8 \times \qty{2.5}{m}$~; son air est
à $P = \qty{1.013e5}{Pa}$ et $T = \qty{293}{K}$, de masse molaire
$M = \qty{29.0}{g/mol}$, avec $79\%$ de N$_2$ et $21\%$ de O$_2$ en
molécules. $R = \qty{8.314}{J/(mol.K)}$,
$k_B = \qty{1.38e-23}{J/K}$, $N_A = \num{6.02e23}$, $g = \qty{9.81}{m/s^2}$.

\medskip
\noindent\textbf{Partie I --- Compter et peser.}
\begin{enumerate}
  \item Quantité de matière (en \omterm{def:b1:kinetic-theory:scales}{moles}) et nombre de molécules dans la salle.
  \item Masse de cet air~; la comparer à celle d'une personne.
  \item Densité particulaire $n^*$ et distance moyenne entre molécules
        ($n^{*-1/3}$)~; comparer au diamètre moléculaire \qty{0.37}{nm}.
  \item \omterm{def:b1:kinetic-theory:kinetictemp}{Vitesse quadratique moyenne} de N$_2$ et de O$_2$~; pourquoi le gaz
        le plus lourd va-t-il moins vite à même \omterm{def:b1:kinetic-theory:temperature}{température}~?
  \item \omterm{thm:b1:work-and-energy:ke}{Énergie cinétique} moyenne d'une molécule~; \omterm{thm:b1:work-and-energy:ke}{énergie cinétique} de
        translation totale de l'air de la salle~; comparer à l'\omterm{thm:b1:work-and-energy:ke}{énergie
        cinétique} d'une voiture à \qty{100}{km/h}.
  \item \omterm{rem:b1:kinetic-theory:validity}{Libre parcours moyen} $\lambda \approx 1/(\sqrt2\pi d^2n^*)$ et
        temps moyen entre deux chocs d'une molécule.
\end{enumerate}

\medskip
\noindent\textbf{Partie II --- La \omterm{def:b1:kinetic-theory:pressure}{pression}, par les chocs.}
\begin{enumerate}[resume]
  \item Avec le modèle à six directions et $u = \qty{500}{m/s}$, calculer
        le nombre de chocs par seconde sur \qty{1}{cm^2} de paroi.
  \item Calculer la quantité de mouvement cédée par seconde, et vérifier
        qu'elle reproduit la \omterm{def:b1:kinetic-theory:pressure}{pression} atmosphérique.
  \item Force totale de l'air sur une paroi de \qty{8}{m} sur \qty{2.5}{m}~;
        pourquoi la paroi ne bouge-t-elle pas~?
  \item La salle est réchauffée à \qty{303}{K}, fenêtres fermées et
        étanches~: nouvelle \omterm{def:b1:kinetic-theory:pressure}{pression}, et force résultante sur la paroi.
  \item Même réchauffement fenêtre ouverte~: qu'est-ce qui sort, et
        combien~?
\end{enumerate}

\medskip
\noindent\textbf{Partie III --- La poussée d'Archimède.}
\begin{enumerate}[resume]
  \item Masse volumique de l'air de la salle, et de l'air à \qty{373}{K}
        sous la même \omterm{def:b1:kinetic-theory:pressure}{pression}.
  \item Une montgolfière de \qty{2800}{m^3} est remplie d'air à
        \qty{373}{K}~: poids de l'air chaud, poids de l'air froid
        déplacé, et force ascensionnelle disponible (Archimède,
        \cref{ch:b1:fluid-statics}).
  \item L'enveloppe, la nacelle, le brûleur et le gaz pèsent
        \qty{300}{kg}~: combien de passagers de \qty{75}{kg} peut-elle
        emporter~?
  \item À quelle \omterm{def:b1:kinetic-theory:temperature}{température} faudrait-il chauffer l'air pour doubler la
        force ascensionnelle, et pourquoi ne le fait-on pas (l'enveloppe
        est en nylon)~?
  \item Un ballon d'hélium de même volume à \qty{293}{K}
        ($M = \qty{4}{g/mol}$)~: quelle force ascensionnelle~? Pourquoi
        les ballons touristiques emploient-ils malgré tout l'air chaud~?
\end{enumerate}

\medskip
\noindent\textbf{Partie IV --- L'air raréfié et la vitesse du son.}
\begin{enumerate}[resume]
  \item \omterm{def:b1:kinetic-theory:pressure}{Pression} partielle de l'oxygène dans la salle.
  \item Au sommet de l'Everest, $P = \qty{0.33}{atm}$ et $T \approx \qty{250}{K}$~:
        densité particulaire de l'oxygène comparée à celle de la salle.
  \item La \omterm{def:b1:kinetic-theory:pressure}{pression} de l'atmosphère décroît avec l'altitude à peu près
        comme $P(z) = P_0\eu^{-z/H}$ avec $H = RT/Mg$ (\cref{ch:b1:fluid-statics})~:
        calculer $H$ pour \qty{260}{K} et l'altitude où $P = P_0/2$.
  \item La vitesse du son dans un \omterm{def:b1:kinetic-theory:model}{gaz parfait} vaut $c_s = \sqrt{\gamma RT/M}$
        avec $\gamma = 1.4$ pour l'air (\cref{ch:b1:first-law})~: la
        calculer à \qty{293}{K} et la comparer à la \omterm{def:b1:kinetic-theory:kinetictemp}{vitesse quadratique
        moyenne}~; interpréter cette proximité.
  \item Comment $c_s$ varie-t-elle avec $T$~? Vitesse du son dans la
        salle à \qty{303}{K}, et dans l'hélium à \qty{293}{K}
        ($\gamma = 5/3$)~: pourquoi une voix à l'hélium sonne-t-elle aiguë~?
  \item Au sommet de l'Everest, le son va-t-il plus vite ou moins vite
        que dans la salle~? De combien~?
  \item On remplace l'air de la salle par de l'argon aux mêmes $P$ et
        $T$~: qu'est-ce qui change dans le décompte, la masse, les
        vitesses, l'\omterm{def:b1:kinetic-theory:U}{énergie interne}~?
  \item \omterm{def:b1:kinetic-theory:U}{Énergie interne} de l'air de la salle ($U = \tfrac52 nRT$) et
        énergie pour la réchauffer de \qty{10}{K} à volume constant.
  \item Récapituler en trois lignes ce qu'a donné le modèle cinétique~:
        quelles grandeurs macroscopiques il a expliquées, et à partir de
        quelle unique moyenne microscopique.
\end{enumerate}
\end{problem}
